\section{A DSL for UI generation}

Inspired by \href{https://www.openui.com/}{OpenUI} and \href{https://gen-ui-lang.github.io/}{GenUI Lang}, we propose a DSL to generate interactive UI elements.\\

\noindent This would require:
\begin{itemize}
    \item Designing a declarative and reactive DSL which should have
    	\begin{itemize}
    		\item Control flow statements 
    		\item Variables and assignments
    		\item Functions and scoping
    		\item State management at runtime
    	\end{itemize}
    \item The implementation will be done in OCAML
    \item If time allows, we aim for a proof-of-concept covering the complete flow of user prompt $\rightarrow$ UI code generation by LLM $\rightarrow$ UI rendering in the browser.
\end{itemize}

The overall goal is, to create a DSL which would allow the programmer to write a UI in one language and be able to dispatch it through multiple different runtimes. The LLM part of the design serves as a constraint for direction: a compact language allows fast generation.

An example program would look as such:
\begin{center}
	\begin{lstlisting}
		fn indexed_card(idx)  {
			card {
				text: "Counter: " ++ idx
			}
		}
		
		i := 5
		row {
			button {
				click: i++
			}
			indexed_card(i)
		}
	\end{lstlisting}
\end{center}

The core elements of the project remain language semantics definition and proof,  and the web + LLM integration, although important, mostly serve as a demo.

The project engages with formal semantics as the DSL and its runtime will have to be properly defined and shown to behave as the design intends.